\section{Introduction} \label{sec:intro}

Security evaluation of AI-based autonomous platforms remains fragmented.
Most published work tests a single attack vector against a single defense
in an ad-hoc simulation environment, reports one or two metrics, and
rarely shares a reusable experimental setup.  As a result, it is
difficult to compare defense strategies across studies, reproduce
published results, or extend evaluations to new attack types without
rebuilding the entire test infrastructure.

Several recent efforts have advanced individual pieces of the puzzle.
CARE~\cite{care2024} provides a comprehensive cross-evaluation of
network intrusion detection systems across 15 attacks, 8 defenses, and
12 detectors, but operates on generic network-traffic datasets with no
simulator and no mission-level context.
DYNAMITE~\cite{dynamite2025} demonstrates that no single defense is
universally effective for IoT traffic and proposes ML-based defense
selection, yet targets flat network telemetry rather than autonomous
system command-and-control streams.
UAVIDS-2025~\cite{uavids2025} contributes a labeled UAV intrusion
detection dataset, but provides only static records---no orchestrator, no
pluggable defenses, and no real-time logging.
MIXED-SENSE~\cite{mixedsense2025} offers a mixed-reality UAV testing
framework with mission-level feedback, but supports only a single attack
type (false data injection).

Table~\ref{tab:comparison} summarizes the gap.  No existing framework
simultaneously supports \emph{multiple} attack types, \emph{multiple}
defense strategies, a \emph{runnable} orchestrator for autonomous/CPS
platforms, and \emph{mission-level} impact metrics under identical
experimental conditions.

\begin{table}[t]
\centering
\caption{Comparison of related evaluation frameworks.}
\label{tab:comparison}
\begin{tabular}{lccccc}
\toprule
\textbf{Framework} & \rotatebox{70}{\textbf{Multi-attack}} & \rotatebox{70}{\textbf{Multi-defense}} & \rotatebox{70}{\textbf{CPS/Autonomous}} & \rotatebox{70}{\textbf{Runnable}} & \rotatebox{70}{\textbf{Mission metrics}} \\
\midrule
CARE~\cite{care2024}            & \checkmark & \checkmark &            & \checkmark &            \\
DYNAMITE~\cite{dynamite2025}    & \checkmark & \checkmark &            &            &            \\
UAVIDS-2025~\cite{uavids2025}   & \checkmark &            & \checkmark &            &            \\
MIXED-SENSE~\cite{mixedsense2025}& &         \checkmark & \checkmark & \checkmark & partial    \\
FIMSim~\cite{fimsim}            &            &            & \checkmark & \checkmark &            \\
\textbf{Ours}                   & \checkmark & \checkmark & \checkmark & \checkmark & \checkmark \\
\bottomrule
\end{tabular}
\end{table}

This paper presents an evaluation harness that fills this gap.
The framework provides:
\begin{itemize}
    \item A pluggable architecture with abstract base classes for attacks,
          defenses, detectors, and platform simulators, allowing new modules
          to be added without modifying the orchestrator.
    \item An orchestrator that runs multiple attack types against multiple
          defense strategies under identical network and timing conditions,
          with synchronized event logging.
    \item Standardized, automatically computed metrics: detection rate,
          false-positive rate, detection latency, block rate, accuracy
          recovery, and mission impact.
    \item Support for both static and randomized (non-stationary) attack
          scheduling to evaluate defense robustness under varying conditions.
\end{itemize}

We use the harness to compare two defense strategies---a Markov
transition-probability filter and an Isolation Forest anomaly
detector with adaptive blocking---against six MAVLink injection attack
types.  We conduct four experiments of increasing difficulty: static
attacks without defense, static attacks with Markov defense, randomized
attacks with Markov defense, and randomized attacks with Isolation Forest
defense.  The results demonstrate that pattern-based defenses degrade
significantly under non-stationary attack traffic, while feature-based ML
defenses maintain effectiveness across attack variations.
